% This file should be rename with the speaker's last name. (e.g : smith)

% Write the author names in the order you wish them to appear
% and underline the one who will give the talk.

% Only the speaker's email address is required.

\begin{abstract_online}{Teaching the residue theorem and its applications with a {\sc Cas}}{%
        \underline{Jos\'e Luis Gal\'an-Garc\'{\i}a}$^{1}$, Gabriel Aguilera-Venegas$^{1}$, Pedro Rodr\'{\i}guez-Cielos$^{1}$, Yolanda Padilla-Dom\'{\i}nguez$^{1}$, Mar\'{\i}a \'Angeles Gal\'an-Garc\'{\i}a$^{1}$}{[jlgalan@uma.es]
        }{%
        $^1$ Applied Mathematics, University of M\'alaga, M\'alaga, Spain% \newline{}$^2$ Physics Department, Lancaster University, Lancaster, UK
}

The residue theorem is one of the most interesting result in Complex Analysis which allows not only computations in $\mathbb{C}$, the Field of Complex Numbers, but also provides many applications in the Field of Real Numbers $\mathbb{R}$.
 
In [1] we introduced the file {\tt RESIDUE.MTH}, developed in the {\sc Cas Derive} which main objective was to provide tools for solving integration problems in Complex Analysis using the residue theorem.

In this talk we present the file {\tt ResidueApplications.dfw}, developed also in {\sc Derive}, aimed to teach the residue applications in Engineering degrees.

{\tt ResidueApplications.dfw} can be used as a tutorial in the teaching and learning process of this topic since it provides the results step by step allowing the students to check their computations when they solve an exercise.

The programs developed in this tutorial can be grouped in the following blocks:
\begin{enumerate}
\item Compute of residues.
\item Compute of complex integrals using the residue theorem.
\item Applications of the residue theorem to compute integrals in $\mathbb{R}$:
\begin{enumerate}
\item Trigonometric integrals.
\item Improper integrals.
\end{enumerate}
\end{enumerate}
In previous ACA conferences we dealt with the application of the residue theorem to compute improper integrals (see [1] and [2]). In this talk, although we will present an overview of the whole tutorial, we will focus mainly in the computation of trigonometric integrals. 


\textbf{Keywords} \newline{} Residue theorem, Trigonometric integrals, Improper integrals, Stepwise tutorial, {\sc Cas}, {\sc Derive}

\textbf{References} \newline{}
[1] \textsc{J.L. Gal\'an-Garc\'{\i}a, M.\'A. Gal\'an-Garc\'{\i}a, Y. Padilla-Dom\'{\i}nguez, P. Rodr\'{\i}guez-Cielos}, 91.	Residue.mth: solving problems of integration using the residue theorem. In \emph{Proceedings of TIME-2004 Symposium: Technology and its Integration in Mathematics Education}, Montreal (Canada), 2004, ISBN 3-901769-59-5. \newline{}
[2] \textsc{G. Aguilera, J.L. Galán, M.A. Galán, Y. Padilla, P. Rodríguez, R. Rodríguez}, 15.	Teaching improper integrals with {\sc Cas}. In \emph{21st Conference on Applications of Computer Algebra - ACA 2015}, 90--91. Kalamata (Greece), 2015, ISBN 978-84-16954-87-2. \newline{}
[3] \textsc{J.L. Gal\'an-Garc\'{\i}a, G. Aguilera-Venegas, P. Rodr\'{\i}guez-Cielos, Y. Padilla-Dom\'{\i}nguez, M.\'A. Gal\'an-Garc\'{\i}a}, 2.	New rules for improving CAS capabilities when computing improper integrals. Applications in Math Education. In \emph{24th Conference on Applications of Computer Algebra - ACA 2018}, Francisco Botana, Felipe Gago and Manuel Ladra (eds.), 63--63. Servizo de Publicaci\'ons e Intercambio Cient\'{\i}fico Campus Vida, Santiago de Compostela (Spain), 2018, ISBN 978-84-16954-87-2. \newline{}
\end{abstract_online}
    